\documentclass[10pt]{article}
\usepackage[margin=0.72in]{geometry}
\usepackage[T1]{fontenc}
\usepackage{amsmath}
\usepackage{enumitem}
\usepackage{hyperref}
\usepackage{xcolor}
\hypersetup{colorlinks=true, urlcolor=blue, linkcolor=black}
\setlist{nosep,leftmargin=1.35em}
\setlength{\parindent}{0pt}
\setlength{\parskip}{4pt}
\begin{document}

\begin{center}
{\Large\bfseries Anonymous Long-Video QA Artifact}\\[3pt]
{\small Technical document for reproducibility and anonymous submission}
\end{center}

\section*{1. Problem and artifact scope}

The task is open-ended question answering over short films. Each example
contains a movie identifier, a natural-language question, and an answer that
must be submitted without exposing evaluation labels. The central engineering
problem is long-range evidence integration: a question may depend on a name
introduced early in the film, a relation that changes later, or the final state
of an event chain.

This artifact releases inference and submission code only. Raw videos,
subtitles, model weights, generated predictions, credentials, and evaluation
records are external inputs. The boundary supports anonymous review and
prevents accidental disclosure of test information.

\section*{2. Movie-level answer-bank construction}

The main runner groups questions by \texttt{video\_id}. For each movie it
builds one request containing a timestamped transcript and chronological visual
observations. All questions for that movie are placed in the same request. The
model returns exactly one JSON object whose keys are the original question IDs
and whose values are concise English answers.

The computation is:
\[
  (E_m,Q_m) \longrightarrow
  \{a_{m,1},a_{m,2},\ldots,a_{m,n_m}\},
\]
where $E_m$ is shared evidence for movie $m$, $Q_m$ is its complete question
set, and the answer bank is emitted in one autoregressive context. The prompt
asks for consistent names, aliases, relationships, chronology, causes,
repeated objects, and ending state. It forbids facts unsupported by the
transcript or frames.

The evidence contract is explicit. Transcript cues are serialized with start
timestamps. Frames are sampled chronologically, resized to a fixed maximum
side, encoded as JPEG, and placed after the text prompt with timestamp markers.
Long transcripts use a head-and-tail bound so setup and ending evidence remain
visible.

\newpage

\section*{3. Reproducibility and failure handling}

The runner accepts an OpenAI-compatible Responses endpoint and reads its API
key from a process environment variable. The key is never written to output
files or logs. Movie requests can be retried for transient provider failures;
answers are persisted only after the complete JSON object has been parsed and
validated against every question ID in that movie. Missing credentials,
malformed JSON, missing answers, and incomplete movie banks are errors.

The all-or-nothing movie boundary prevents partial requests from silently
mixing answers produced under different evidence or decoding conditions. A
separate conversion script maps a validated answer bank into the competition
schema. The audit script checks exact ID coverage, duplicate rows, empty
predictions, and the required header.

\section*{4. Competition-facing compliance}

The release follows the operational conventions used by the SLoMO/EvalAI
workflow:
\begin{itemize}
  \item use the live challenge phase and current metadata fields;
  \item submit exactly \texttt{question\_id,prediction};
  \item include every required ID exactly once with a non-empty prediction;
  \item keep public and private annotations separate;
  \item do not use answer keys, options, private labels, or leaderboard output
        during inference.
\end{itemize}

The repository starts from a fresh anonymous Git history. Identity-bearing
material is outside the tree: account names, team names, local absolute paths,
API tokens, private data, generated answer banks, and submission IDs. The live
competition page and EvalAI portal remain authoritative for deadlines,
availability, metadata, and upload policy:
\begin{center}
\url{https://slomo-workshop.github.io/eccv2026/}\\
\url{https://eval.ai/}
\end{center}

\section*{5. Known limitations}

The method depends on transcript quality, frame coverage, model endpoint
availability, and context budget. This artifact does not claim that every
endpoint or model identifier remains available. It does not package
competition data or an official score. Reproduction requires authorized data
access and an independent local audit of the generated CSV.

\end{document}
